% ============================================================================
\section{Pipelining}
\label{sec:pipe}

Everything on pipelining: the idea, the 5-stage MIPS pipeline and its timing,
structural, data and control hazards, forwarding, compiler scheduling and branch
schemes.

% ----------------------------------------------------------------------------
\subsection{What to know}

\subsubsection{The idea}
\begin{itemize}[leftmargin=*]
  \item \textbf{Pipelining} overlaps the execution of several instructions, each in
  a different stage. It improves \textbf{throughput}, not the latency of one
  instruction (latency even grows a little from pipeline-register overhead).
  \item The pipeline rate is limited by the \textbf{slowest stage}; unbalanced
  stages and the time to fill and drain reduce the speedup. Ideal speedup $=$
  number of stages.
  \item Laundry: wash 30, dry 40, fold 20 min, 4 loads: sequential
  $4\times90=360$ min, pipelined $30+4\times40+20=210$ min.
\end{itemize}

\subsubsection{The 5-stage MIPS pipeline}
\begin{center}
\small
\begin{tabular}{@{}lll@{}}
\toprule
Stage & Work & Resource \\
\midrule
IF & fetch instruction at PC, PC $\leftarrow$ PC$+4$ & instruction memory \\
ID & decode, read registers, sign-extend immediate & register file (read) \\
EX & ALU operation or effective address & ALU \\
MEM & load or store & data memory \\
WB & write the result register & register file (write) \\
\bottomrule
\end{tabular}
\end{center}
Pipeline registers IF/ID, ID/EX, EX/MEM, MEM/WB hold each instruction's values
between stages. The register file is \textbf{written in the first half} and
\textbf{read in the second half} of a cycle. MIPS is easy to pipeline: all
instructions are 32 bits, few formats with fixed register fields, memory operands
only in loads and stores, aligned memory accesses.

\subsubsection{Timing}
\begin{itemize}[leftmargin=*]
  \item \textbf{Clocked pipeline}, $k$ stages, $n$ instructions, $s$ stall cycles:
  \[
    \text{Cycles}=k+(n-1)+s,\qquad t_{clk}=\max_i t_i+t_{latch},\qquad
    \text{Time}=\text{Cycles}\times t_{clk}.
  \]
  \item Non-pipelined time per instruction $=\sum_i t_i$ (single-cycle: the clock
  is the longest instruction).
  \item Speedup, balanced stages: $\dfrac{nk}{k+n-1}\to k$. With stalls:
  \[
    \CPI_{pipe}=1+\text{stall cycles per instruction},
  \]
  \[
    \text{Speedup}=\frac{\CPI_{unpipe}}{\CPI_{pipe}}\times\frac{t_{clk,unpipe}}{t_{clk,pipe}}
    =\frac{\text{depth}}{1+\text{stalls/instr}}.
  \]
  \item \textbf{Stages that hand over as soon as the next one is free} (the
  lecturer's figure, no common clock): the first instruction takes $\sum t_i$, each
  later one finishes $\max t_i$ later:
  $T(n)=\sum_i t_i+(n-1)\max_i t_i$. An instruction waits in front of the slowest
  stage for (slowest $-$ slowest earlier stage) longer than the instruction before it.
\end{itemize}

\subsubsection{Hazards}
\begin{definition}[Hazards]
A \textbf{hazard} prevents the next instruction from executing in its cycle.
\textbf{Structural}: two instructions need the same hardware at once.
\textbf{Data}: an instruction needs a result that is not ready.
\textbf{Control}: the next PC depends on a branch not yet resolved.
\end{definition}
\begin{itemize}[leftmargin=*]
  \item \textbf{Structural}: e.g.\ one memory port for instructions and data, so a
  load stalls a later fetch. Fix: split I/D memories (caches) or duplicate the
  resource. It may be accepted when rare and costly to remove.
  \item \textbf{Data dependences}, $i$ before $j$:
  \textbf{RAW} ($j$ reads what $i$ writes, true dependence),
  \textbf{WAR} ($j$ writes what $i$ reads, anti-dependence),
  \textbf{WAW} ($j$ writes what $i$ writes, output dependence).
  \textbf{RAR is not a hazard} (nothing changes). In the in-order 5-stage pipeline
  only RAW causes stalls: reads (ID) are early and writes (WB) late and in order.
  WAR and WAW matter in out-of-order machines (Section~\ref{sec:mp}).
\end{itemize}

\begin{center}
\small
\begin{tabular}{@{}p{5cm}lp{6.2cm}@{}}
\toprule
Producer $\to$ consumer at distance $d$ & Stalls & Why \\
\midrule
No forwarding & $\max(0,3-d)$ & consumer's ID must be in the same cycle as producer's WB or later \\
Forwarding, ALU $\to$ ALU & 0 & EX/MEM $\to$ EX ($d=1$), MEM/WB $\to$ EX ($d=2$) \\
Forwarding, load $\to$ use, $d=1$ & 1 & data ready only at the end of MEM \\
Forwarding, load $\to$ use, $d\ge2$ & 0 & MEM/WB $\to$ EX \\
Forwarding, load $\to$ store data & 0 & MEM/WB $\to$ MEM \\
\bottomrule
\end{tabular}
\end{center}
\begin{itemize}[leftmargin=*]
  \item \textbf{Forwarding} (bypassing, hardware): feed the ALU from the pipeline
  registers instead of waiting for WB. The \textbf{load-use} hazard remains:
  hardware \textbf{interlocks} detect it and stall 1 cycle.
  \item \textbf{Compiler scheduling}: reorder independent instructions so no load is
  followed directly by its use. Reordering must keep every RAW, WAR and WAW order;
  if it cannot, rename registers.
  \item \textbf{Drawing a diagram}: a stalled instruction waits in ID
  (\textit{st} after ID); the one behind waits in IF (\textit{st} after IF); the
  next instruction is fetched when the IF stage is free. Total cycles $=$ cycle of
  the last WB.
\end{itemize}

\subsubsection{Control hazards}
\begin{itemize}[leftmargin=*]
  \item Branch penalty $=$ (stage that resolves the branch) $-1$: resolved in MEM
  $\Rightarrow$ 3 cycles; with the zero test and target adder moved into ID
  $\Rightarrow$ 1 cycle.
  \item $\CPI=1+f_{branch}\times(\text{fraction penalized})\times\text{penalty}$.
  \item Four alternatives: (1) \textbf{stall/flush} until resolved; (2)
  \textbf{predict not taken}: keep fetching PC$+4$, squash if taken (only taken
  branches pay); (3) \textbf{predict taken}: no gain in the 5-stage MIPS since the
  target is known no earlier than the outcome; (4) \textbf{delayed branch}: the
  instruction in the delay slot always executes.
  \item Filling the delay slot: (A) from before the branch (best, always useful,
  the branch must not depend on it); (B) from the target (useful when taken; must
  be safe when not); (C) from the fall-through (useful when not taken; must be safe
  when taken). Compilers fill about 60\% of slots, about 80\% of those usefully.
  \item \textbf{Compiler techniques}: instruction scheduling, delay-slot filling,
  static branch prediction, loop unrolling. \textbf{Hardware techniques}:
  forwarding, interlocks, early branch resolution, dynamic branch prediction.
\end{itemize}

% ----------------------------------------------------------------------------
\subsection{Common traps}
\begin{enumerate}[leftmargin=*]
  \item $n$ instructions take $k+n-1$ cycles, not $kn$ and not $n$.
  \item The clock is the \emph{slowest} stage (plus latch overhead), not the average.
  \item Without forwarding a dependence at distance 1 costs 2 stalls (not 3),
  because of the split-cycle register file.
  \item With forwarding only load-use at distance 1 stalls.
  \item Stalls push every later instruction right: the instruction behind a stall
  waits in IF.
  \item When reordering, check WAR and WAW too: moving a load above an instruction
  that still reads the old register value is wrong.
  \item Forwarding is hardware; scheduling and delay slots are compiler techniques.
\end{enumerate}

% ----------------------------------------------------------------------------
\subsection{Lecturer's review questions}

\begin{exercise}[Lecturer 1]
Choose the FALSE statement about techniques to reduce hazards \emph{using
compilers}: (a) forwarding; (b) scheduling; (c) branch prediction; (d) branch delay
slots.
\end{exercise}
\begin{solution}
\textbf{(a)}: forwarding is a hardware datapath feature. Scheduling and delay-slot
filling are done by the compiler, and branch prediction can be static (the compiler
or ISA chooses the predicted direction).
\end{solution}

\begin{exercise}[Lecturer 2]
Which data dependence is not a data hazard? (a) RAR; (b) RAW; (c) WAR; (d) WAW.
\end{exercise}
\begin{solution}
\textbf{(a)}: two reads of the same register can happen in any order.
\end{solution}

\begin{exercise}[Lecturer 3]
How long do 100 tasks take on an ideal 6-stage pipeline with 2~ns per stage?
(a) 12~ns; (b) 200~ns; (c) 210~ns; (d) 212~ns.
\end{exercise}
\begin{solution}
\textbf{(c)}: $(6+100-1)\times2=105\times2=210$~ns. (12~ns is the latency of one
task; 200~ns ignores the fill.)
\end{solution}

\begin{exercise}[Lecturer 4 and 5]
A 5-stage pipeline has stage times IF 5, ID 4, EX 5, MEM 10, WB 4~ns.
(4) The total execution time of 3 consecutive instructions is: (a) 38; (b) 43;
(c) 48; (d) 53~ns.
(5) The stall time of the 4th instruction is: (a) 5; (b) 10; (c) 15~ns; (d) none of
the above.
\end{exercise}
\begin{solution}
The answer choices fit the model in which an instruction moves on as soon as the
next stage is free (no common clock).
\textbf{(4) (c) 48~ns}: the first instruction takes $5+4+5+10+4=28$~ns; MEM, the
slowest stage, accepts one instruction every 10~ns, so each later instruction
finishes 10~ns later: $28+2\times10=48$~ns.
\begin{center}
\small
\begin{tabular}{@{}lccccc@{}}
\toprule
Instr. & IF & ID & EX & wait & MEM, WB \\
\midrule
1 & 0--5 & 5--9 & 9--14 & -- & 14--24, 24--28 \\
2 & 5--10 & 10--14 & 14--19 & 5 & 24--34, 34--38 \\
3 & 10--15 & 15--19 & 19--24 & 10 & 34--44, 44--48 \\
4 & 15--20 & 20--24 & 24--29 & 15 & 44--54, 54--58 \\
\bottomrule
\end{tabular}
\end{center}
\textbf{(5) (c) 15~ns}: instructions arrive in front of MEM every 5~ns (the slowest
earlier stage) but MEM takes 10~ns, so each one waits 5~ns longer than the one
before: $0, 5, 10, 15$~ns.

With a common clock instead, the cycle would be 10~ns and 3 instructions would take
$(5+2)\times10=70$~ns, which is not among the choices.
\end{solution}

\begin{exercise}[Lecturer 6: pipeline diagrams]\label{l:pipe-code}
The code runs on the 5-stage pipeline IF, ID, EX, MEM, WB:
\begin{lstlisting}
(1) lw   $t0, 0($t2)
(2) lw   $t1, 4($t4)
(3) add  $t3, $t0, $t1
(4) sw   $t3, 8($t2)
(5) lw   $t0, 12($t4)
(6) add  $t3, $t0, $t1
(7) sw   $t3, 16($t2)
\end{lstlisting}
\begin{enumerate}[label=(\alph*),leftmargin=*]
  \item Without forwarding, fill in the pipeline table, writing STALL where needed.
  \item Repeat with forwarding.
  \item Reorder the instructions so they run on the forwarding pipeline with as few
  stalls as possible.
\end{enumerate}
\end{exercise}
\begin{solution}
Dependences: RAW (1)$\to$(3) \texttt{\$t0}, (2)$\to$(3) \texttt{\$t1},
(3)$\to$(4) \texttt{\$t3}, (5)$\to$(6) \texttt{\$t0}, (2)$\to$(6) \texttt{\$t1},
(6)$\to$(7) \texttt{\$t3}; WAR (3)$\to$(5) on \texttt{\$t0}, (4)$\to$(6) on
\texttt{\$t3}; WAW (1)$\to$(5), (3)$\to$(6). \textit{st} $=$ STALL.

\textbf{(a) No forwarding}: a consumer's ID (register read) must be in the cycle of
the producer's WB or later.
\begin{center}\scriptsize\setlength{\tabcolsep}{1.1pt}
\begin{tabular}{@{}lccccccccccccccccccc@{}}
\toprule
Instruction & 1 & 2 & 3 & 4 & 5 & 6 & 7 & 8 & 9 & 10 & 11 & 12 & 13 & 14 & 15 & 16 & 17 & 18 & 19 \\
\midrule
(1) \texttt{lw \$t0,0(\$t2)} & IF & ID & EX & MEM & WB &  &  &  &  &  &  &  &  &  &  &  &  &  &  \\
(2) \texttt{lw \$t1,4(\$t4)} &  & IF & ID & EX & MEM & WB &  &  &  &  &  &  &  &  &  &  &  &  &  \\
(3) \texttt{add \$t3,\$t0,\$t1} &  &  & IF & ID & \textit{st} & \textit{st} & EX & MEM & WB &  &  &  &  &  &  &  &  &  &  \\
(4) \texttt{sw \$t3,8(\$t2)} &  &  &  & IF & \textit{st} & \textit{st} & ID & \textit{st} & \textit{st} & EX & MEM & WB &  &  &  &  &  &  &  \\
(5) \texttt{lw \$t0,12(\$t4)} &  &  &  &  &  &  & IF & \textit{st} & \textit{st} & ID & EX & MEM & WB &  &  &  &  &  &  \\
(6) \texttt{add \$t3,\$t0,\$t1} &  &  &  &  &  &  &  &  &  & IF & ID & \textit{st} & \textit{st} & EX & MEM & WB &  &  &  \\
(7) \texttt{sw \$t3,16(\$t2)} &  &  &  &  &  &  &  &  &  &  & IF & \textit{st} & \textit{st} & ID & \textit{st} & \textit{st} & EX & MEM & WB \\
\bottomrule
\end{tabular}
\end{center}
(3) waits for \texttt{\$t1} (WB in 6); (4) waits for \texttt{\$t3} (WB in 9); (6)
waits for \texttt{\$t0} (WB in 13); (7) waits for \texttt{\$t3} (WB in 16).
\textbf{19 cycles, 8 stall cycles.}

\textbf{(b) Forwarding}: only the two load-use pairs (2)$\to$(3) and (5)$\to$(6)
stall, 1 cycle each; (3)$\to$(4) and (6)$\to$(7) are forwarded from EX/MEM.
\begin{center}\scriptsize\setlength{\tabcolsep}{1.6pt}
\begin{tabular}{@{}lccccccccccccc@{}}
\toprule
Instruction & 1 & 2 & 3 & 4 & 5 & 6 & 7 & 8 & 9 & 10 & 11 & 12 & 13 \\
\midrule
(1) \texttt{lw \$t0,0(\$t2)} & IF & ID & EX & MEM & WB &  &  &  &  &  &  &  &  \\
(2) \texttt{lw \$t1,4(\$t4)} &  & IF & ID & EX & MEM & WB &  &  &  &  &  &  &  \\
(3) \texttt{add \$t3,\$t0,\$t1} &  &  & IF & ID & \textit{st} & EX & MEM & WB &  &  &  &  &  \\
(4) \texttt{sw \$t3,8(\$t2)} &  &  &  & IF & \textit{st} & ID & EX & MEM & WB &  &  &  &  \\
(5) \texttt{lw \$t0,12(\$t4)} &  &  &  &  &  & IF & ID & EX & MEM & WB &  &  &  \\
(6) \texttt{add \$t3,\$t0,\$t1} &  &  &  &  &  &  & IF & ID & \textit{st} & EX & MEM & WB &  \\
(7) \texttt{sw \$t3,16(\$t2)} &  &  &  &  &  &  &  & IF & \textit{st} & ID & EX & MEM & WB \\
\bottomrule
\end{tabular}
\end{center}
\textbf{13 cycles, 2 stall cycles.}

\textbf{(c) Reordering.} Moving (5) up between (2) and (3) looks attractive, but
(5) overwrites \texttt{\$t0} before (3) reads the value loaded by (1): a WAR
violation that changes the result. Keeping all dependences, (3) must follow both
(1) and (2), and nothing independent can go between them, so one load-use stall
remains. Swapping (4) and (5) removes the second one:
\begin{center}\scriptsize\setlength{\tabcolsep}{1.6pt}
\begin{tabular}{@{}lcccccccccccc@{}}
\toprule
Instruction & 1 & 2 & 3 & 4 & 5 & 6 & 7 & 8 & 9 & 10 & 11 & 12 \\
\midrule
(1) \texttt{lw \$t0,0(\$t2)} & IF & ID & EX & MEM & WB &  &  &  &  &  &  &  \\
(2) \texttt{lw \$t1,4(\$t4)} &  & IF & ID & EX & MEM & WB &  &  &  &  &  &  \\
(3) \texttt{add \$t3,\$t0,\$t1} &  &  & IF & ID & \textit{st} & EX & MEM & WB &  &  &  &  \\
(5) \texttt{lw \$t0,12(\$t4)} &  &  &  & IF & \textit{st} & ID & EX & MEM & WB &  &  &  \\
(4) \texttt{sw \$t3,8(\$t2)} &  &  &  &  &  & IF & ID & EX & MEM & WB &  &  \\
(6) \texttt{add \$t3,\$t0,\$t1} &  &  &  &  &  &  & IF & ID & EX & MEM & WB &  \\
(7) \texttt{sw \$t3,16(\$t2)} &  &  &  &  &  &  &  & IF & ID & EX & MEM & WB \\
\bottomrule
\end{tabular}
\end{center}
Order (1) (2) (3) (5) (4) (6) (7): \textbf{12 cycles, 1 stall}. If a free register
may be used, rename \texttt{\$t0} in (5)--(6) to \texttt{\$t5} and \texttt{\$t3} in
(6)--(7) to \texttt{\$t6}; then (5) can move up and there are no stalls:
\begin{center}\scriptsize\setlength{\tabcolsep}{1.6pt}
\begin{tabular}{@{}lccccccccccc@{}}
\toprule
Instruction & 1 & 2 & 3 & 4 & 5 & 6 & 7 & 8 & 9 & 10 & 11 \\
\midrule
(1) \texttt{lw \$t0,0(\$t2)} & IF & ID & EX & MEM & WB &  &  &  &  &  &  \\
(2) \texttt{lw \$t1,4(\$t4)} &  & IF & ID & EX & MEM & WB &  &  &  &  &  \\
(5) \texttt{lw \$t5,12(\$t4)} &  &  & IF & ID & EX & MEM & WB &  &  &  &  \\
(3) \texttt{add \$t3,\$t0,\$t1} &  &  &  & IF & ID & EX & MEM & WB &  &  &  \\
(4) \texttt{sw \$t3,8(\$t2)} &  &  &  &  & IF & ID & EX & MEM & WB &  &  \\
(6) \texttt{add \$t6,\$t5,\$t1} &  &  &  &  &  & IF & ID & EX & MEM & WB &  \\
(7) \texttt{sw \$t6,16(\$t2)} &  &  &  &  &  &  & IF & ID & EX & MEM & WB \\
\bottomrule
\end{tabular}
\end{center}
\textbf{11 cycles, 0 stalls} ($=5+7-1$, the ideal).
\end{solution}

% ----------------------------------------------------------------------------
\subsection{Exam-style questions}
Solutions are in Section~\ref{sec:pipe-sol}.

\begin{exercise}[Clocked pipeline]\label{q:pipe-clock}
Use the stage times of Lecturer question 4 (5, 4, 5, 10, 4~ns) and a 1~ns
pipeline-register overhead.
\begin{enumerate}[label=(\alph*),leftmargin=*]
  \item Find the clock cycle and the time for 1000 instructions on the clocked
  pipeline, and on a non-pipelined machine that takes 28~ns per instruction.
  \item MEM is split into two 5~ns stages. Repeat (a) and give the new speedup.
  \item Why is the speedup in (a) far below 5?
\end{enumerate}
\end{exercise}

\begin{exercise}[Finding hazards]\label{q:pipe-hz}
\leavevmode
\begin{lstlisting}
(1) add $t1, $t2, $t3
(2) sub $t4, $t1, $t5
(3) lw  $t2, 0($t4)
(4) or  $t1, $t6, $t7
(5) sw  $t2, 4($t1)
\end{lstlisting}
\begin{enumerate}[label=(\alph*),leftmargin=*]
  \item List every RAW, WAR and WAW dependence.
  \item Which of them can cause stalls in the 5-stage pipeline? Why not the others?
  \item Draw the diagram without forwarding and with forwarding. Give the cycles.
\end{enumerate}
\end{exercise}

\begin{exercise}[Branch schemes]\label{q:pipe-br}
20\% of instructions are branches and 60\% of them are taken. All other
instructions have $\CPI=1$. Find the CPI for:
(a) stall, branch resolved in MEM;
(b) predict not taken, resolved in MEM;
(c) stall, resolved in ID;
(d) predict not taken, resolved in ID;
(e) one delay slot (resolved in ID) filled usefully 50\% of the time.
(f) Give the speedup of (e) over (a).
\end{exercise}

\begin{exercise}[Delay slots]\label{q:pipe-slot}
\leavevmode
\begin{lstlisting}
      add  $s1, $s2, $s3
      beq  $s4, $zero, L
      (delay slot)
      sub  $t0, $t1, $t2
      ...
L:    or   $t3, $t4, $t5
\end{lstlisting}
\begin{enumerate}[label=(\alph*),leftmargin=*]
  \item Which instruction should fill the slot, and why is it the best choice?
  \item If the branch were \texttt{beq \$s1,\$zero,L}, why could \texttt{add} not
  be used? What are the other two options and their conditions?
\end{enumerate}
\end{exercise}

% ----------------------------------------------------------------------------
\subsection{Solutions}
\label{sec:pipe-sol}

\begin{noot}{Exercise~\ref{q:pipe-clock}}
\begin{enumerate}[label=(\alph*),leftmargin=*]
  \item $t_{clk}=10+1=11$~ns. Pipelined: $(5+999)\times11=11{,}044$~ns.
  Non-pipelined: $1000\times28=28{,}000$~ns. Speedup $=2.54$.
  \item 6 stages, $t_{clk}=5+1=6$~ns: $(6+999)\times6=6{,}030$~ns, speedup
  $28{,}000/6{,}030=4.64$.
  \item The stages are unbalanced: the 10~ns MEM stage sets the clock, so the
  4--5~ns stages idle half of each cycle, and the latch overhead adds to every
  stage.
\end{enumerate}
\end{noot}

\begin{noot}{Exercise~\ref{q:pipe-hz}}
\begin{enumerate}[label=(\alph*),leftmargin=*]
  \item RAW: (1)$\to$(2) \texttt{\$t1}, (2)$\to$(3) \texttt{\$t4}, (3)$\to$(5)
  \texttt{\$t2}, (4)$\to$(5) \texttt{\$t1}. WAR: (1)$\to$(3) \texttt{\$t2},
  (2)$\to$(4) \texttt{\$t1}. WAW: (1)$\to$(4) \texttt{\$t1}.
  \item Only RAW: registers are read in ID and written in WB, in program order, so
  a later instruction can never write before an earlier one reads (WAR) or writes
  (WAW).
  \item Without forwarding (15 cycles, 6 stalls):
\begin{center}\scriptsize\setlength{\tabcolsep}{1.6pt}
\begin{tabular}{@{}lccccccccccccccc@{}}
\toprule
Instruction & 1 & 2 & 3 & 4 & 5 & 6 & 7 & 8 & 9 & 10 & 11 & 12 & 13 & 14 & 15 \\
\midrule
(1) \texttt{add \$t1,\$t2,\$t3} & IF & ID & EX & MEM & WB &  &  &  &  &  &  &  &  &  &  \\
(2) \texttt{sub \$t4,\$t1,\$t5} &  & IF & ID & \textit{st} & \textit{st} & EX & MEM & WB &  &  &  &  &  &  &  \\
(3) \texttt{lw \$t2,0(\$t4)} &  &  & IF & \textit{st} & \textit{st} & ID & \textit{st} & \textit{st} & EX & MEM & WB &  &  &  &  \\
(4) \texttt{or \$t1,\$t6,\$t7} &  &  &  &  &  & IF & \textit{st} & \textit{st} & ID & EX & MEM & WB &  &  &  \\
(5) \texttt{sw \$t2,4(\$t1)} &  &  &  &  &  &  &  &  & IF & ID & \textit{st} & \textit{st} & EX & MEM & WB \\
\bottomrule
\end{tabular}
\end{center}
  With forwarding no stall is needed: (2), (3) and (5) get \texttt{\$t1},
  \texttt{\$t4} and \texttt{\$t1} from EX/MEM, and the load (3) is two
  instructions before the store that uses \texttt{\$t2} (9 cycles):
\begin{center}\scriptsize\setlength{\tabcolsep}{1.6pt}
\begin{tabular}{@{}lccccccccc@{}}
\toprule
Instruction & 1 & 2 & 3 & 4 & 5 & 6 & 7 & 8 & 9 \\
\midrule
(1) \texttt{add \$t1,\$t2,\$t3} & IF & ID & EX & MEM & WB &  &  &  &  \\
(2) \texttt{sub \$t4,\$t1,\$t5} &  & IF & ID & EX & MEM & WB &  &  &  \\
(3) \texttt{lw \$t2,0(\$t4)} &  &  & IF & ID & EX & MEM & WB &  &  \\
(4) \texttt{or \$t1,\$t6,\$t7} &  &  &  & IF & ID & EX & MEM & WB &  \\
(5) \texttt{sw \$t2,4(\$t1)} &  &  &  &  & IF & ID & EX & MEM & WB \\
\bottomrule
\end{tabular}
\end{center}
\end{enumerate}
\end{noot}

\begin{noot}{Exercise~\ref{q:pipe-br}}
(a) $1+0.2\times3=1.6$. (b) $1+0.2\times0.6\times3=1.36$. (c) $1+0.2\times1=1.2$.
(d) $1+0.2\times0.6\times1=1.12$. (e) the slot is wasted half the time:
$1+0.2\times0.5\times1=1.1$. (f) $1.6/1.1=1.45$.
\end{noot}

\begin{noot}{Exercise~\ref{q:pipe-slot}}
\begin{enumerate}[label=(\alph*),leftmargin=*]
  \item \texttt{add \$s1,\$s2,\$s3} from before the branch (strategy A): the branch
  does not use \texttt{\$s1}, and the instruction must execute on both paths, so
  the slot is always useful.
  \item The branch would need the \texttt{add}'s result, so it cannot be moved past
  it. From the target (B): copy \texttt{or} into the slot; useful when taken, and
  \texttt{\$t3} must be unused (dead) on the fall-through path. From the
  fall-through (C): move \texttt{sub} into the slot; useful when not taken, and
  \texttt{\$t0} must be dead at \texttt{L}.
\end{enumerate}
\end{noot}

